\scanpage{47}

% B01 | head, underlined proposition
\uline{Proposition} Soit $k$ un corps, $n$ un entier. Conditions équivalentes

% B02 | upper, clauses a and b
a) $H^i(k,G)=0$ pour $i>n$ et tout groupe alg. comm. $G$ déf$/k$

% B03 | upper, clause b
b) $H^i(k,\Gamma\deleted{\unsure{G|0}})=0$ \inserted{\uline{fini} $\Gamma$ déf$/k$}

% B04 | upper, clause b prime, qualification above
b') $H^i(k,\Gamma)=0$ pour $i>n$ et tout \deleted{groupe espace vectoriel} \inserted{groupe \inserted{fini} $\Gamma$ qui est un} esp. vectoriel sur un corps $Z/\ell$, dans lequel $\mathfrak g_k$ opère (avec noyau ouvert) irréductiblement

% B05 | upper left, oblique at b and b prime
% slanted
\margin{$b,b'$ sont tjrs vérifiés $(\unsure{n\geq1|n>1})$ si $\Gamma$ est un $p$-groupe \deleted{\illegible{1}}}

% B06 | upper middle, clause c, equals with no right operand
c) $H^i(k,G)=$ pour $i>n$ et tout groupe alg. comm. $G$ déf$/k$ \deleted{connexe et affine} devenant isomorphe à un groupe \deleted{\unsure{tore|torus}} $(G_m)^n$ sur la clôture galoisienne de $k$

% B07 | upper middle, clause c prime, left blank and insertion
c') $H^i(k,G)$ \deleted{\unsure{pr|que}} pour $i>n$, chaque fois que $G$ est le groupe des éléments inversibles d'une algèbre \inserted{comm.} \deleted{semi-simple} \inserted{séparable} déf$/k$

% B08 | just below fold
\uline{Démonstration}

% B09 | middle, four branches of implication
\diagram{\begin{tikzcd}[ampersand replacement=\&,row sep=small]
\&a\arrow[dl,Rightarrow]\arrow[dr,Rightarrow]\&\\
 b\arrow[d,Rightarrow]\&\&c\arrow[d,phantom,"\deleted{\Downarrow}" description]\\
 b'\&\&\unsure{a'|c'}
\end{tikzcd}}

% B10 | middle, proof opening
est trivial. Nous prouverons $c\Rightarrow b'$ et $b'\Rightarrow\deleted{b} a$ ce qui achèvera la démonstration

% B11 | below, first reduction in brackets
$b'\Rightarrow b)$ \deleted{En effet $b'\Rightarrow b$} comme on voit par un dévissage immédiat [\deleted{supposant} \inserted{se ramenant} d'abord à $\Gamma$ \deleted{d'ordre} \inserted{annulé par} \inserted{un} $\ell^k$, puis en filtrant par $\Gamma\supset\ell\Gamma\supset\ell^2\Gamma\supset\ldots\supset\ell^k\Gamma=0$, au cas où $\Gamma$ est annulé par $\ell$ i.e. est un vectoriel sur $Z/\ell$, puis enfin au cas irréductible].

% B12 | bottom, fresh implication and struck start
$b\Rightarrow a$ \deleted{En effet considérons le plus grand sous-groupe linéaire $G_1$ de $G$, (qui est défini sur $k$). Alors \illegible{2}} Supposons d'abord que le corps \deleted{soit \`a} $k$ soit \uline{parfait} Soit alors $G_1$ le plus grand sous-groupe linéaire de $G$, on sait qu'il est \uline{défini sur $k$}, et on a la suite exacte

% B13 | foot, exact sequence
\[0\longrightarrow G_1\longrightarrow G\deleted{\unsure{G_1}}\longrightarrow G/G_1\longrightarrow0\qquad(\text{déf}/k)\]

% archive sheet number: 46

